% Tổng quan luồng xử lý SciMKG theo năm ô (a)-(e), nhóm vẽ lại theo cấu trúc tổng quan của bài.
% Khổ 2 cột (figure*), rộng khoảng 17,6 cm. Đơn vị tọa độ là cm, trục y âm hướng xuống.
% Thứ tự đọc: (a) -> (b) ở hàng trên, rồi xuống (c), sang trái qua (d) tới (e).
\begin{tikzpicture}[
  font=\footnotesize, >={Stealth[length=1.7mm]},
  panel/.style={draw=black!40, rounded corners=3pt, fill=black!3},
  plab/.style={font=\footnotesize\bfseries, anchor=north west, inner sep=2pt},
  box/.style={draw, rounded corners=2pt, align=center, inner sep=2pt, font=\scriptsize},
  src/.style={box, fill=yellow!15, minimum height=7.5mm, text width=16mm},
  pre/.style={box, fill=yellow!15, text width=38mm, minimum height=12mm},
  ev/.style={box, fill=orange!15, draw=orange!70!black, text width=22.5mm, minimum height=28.5mm},
  st/.style={box, fill=green!12, draw=green!50!black},
  al/.style={box, fill=cyan!12, draw=cyan!50!black, text width=56mm, minimum height=9mm, align=left},
  app/.style={box, fill=violet!10, draw=violet!60!black, text width=35mm, minimum height=6.2mm},
  arr/.style={->, thick, black!70},
]
% ---------- khung các ô
\draw[panel] (0,0) rectangle (4.2,-3.9);
\draw[panel] (4.8,0) rectangle (17.6,-3.9);
\draw[panel] (0,-4.4) rectangle (3.9,-8.3);
\draw[panel] (4.5,-4.4) rectangle (10.9,-8.3);
\draw[panel] (11.5,-4.4) rectangle (17.6,-8.3);
\node[plab] at (0.05,-0.05) {(a) Tài nguyên giáo dục};
\node[plab] at (4.85,-0.05) {(b) Trích xuất khái niệm: chuỗi EVIA (chạy trên từng đoạn)};
\node[plab] at (0.05,-4.45) {(e) Ứng dụng};
\node[plab] at (4.55,-4.45) {(d) Tổ chức và lưu trữ};
\node[plab] at (11.55,-4.45) {(c) Căn chỉnh đa phương thức};

% ---------- (a) tài nguyên đầu vào và tiền xử lý
\node[src] at (1.1,-1.0)  {Chuẩn chương trình};
\node[src] at (3.1,-1.0)  {Video bài giảng};
\node[src] at (1.1,-1.85) {Slide (PDF)};
\node[src] at (3.1,-1.85) {Bài tập};
\node[pre] (pp) at (2.1,-3.15) {\textbf{Tiền xử lý}\\ ASR, Gemini: đoạn $(t_s,p_s)$\\ cắt clip; slide $\rightarrow$ ảnh};
\draw[arr] (2.1,-2.25) -- (pp.north);

% ---------- (b) EVIA
\node[ev] (e) at (6.225,-2.175)  {\textbf{E} Trích xuất\\[1pt] DeepSeek-V3: tìm danh từ, giữ cụm thuộc môn, in danh sách};
\node[ev] (v) at (8.975,-2.175)  {\textbf{V} Kiểm chứng\\[1pt] Self-Refine: phản hồi từ mơ hồ, rồi sửa; 3 LLM $\times$ 5 vòng};
\node[ev] (i) at (11.725,-2.175) {\textbf{I} Tích hợp\\[1pt] bỏ phiếu: giữ khái niệm có tỉ lệ phiếu $\ge\alpha{=}0{,}7$};
\node[ev] (a) at (14.475,-2.175) {\textbf{A} Bổ sung\\[1pt] ConceptNet (8 quan hệ) và Wikipedia: khái niệm mới, giải thích $\le$20 từ};
\node[st, text width=11mm, minimum height=13mm] (C) at (16.8,-2.175) {tập $C$\\ kèm mốc $t_s$};
\draw[arr] (pp.east) -- (pp.east -| e.west);
\draw[arr] (e) -- (v);
\draw[arr] (v) -- (i);
\draw[arr] (i) -- (a);
\draw[arr] (a) -- (C);

% ---------- (c) căn chỉnh
\node[al] (ai) at (14.55,-5.6) {\textbf{Ảnh} (ngữ nghĩa): GPT-4o $\times$3 trên mỗi ảnh slide, lấy hợp; ưu tiên C4 $>$ C3 $>$ C2 $>$ C1};
\node[al] (av) at (14.55,-6.6) {\textbf{Video} (cấu trúc): khái niệm gắn với clip cắt đúng mốc $t_s$ của đoạn sinh ra nó};
\node[al] (aa) at (14.55,-7.6) {\textbf{Âm thanh}: lời giải thích được đọc bằng TTS thành tệp MP3 cho mỗi khái niệm};
\draw[arr] (C.south) -- (C.south |- ai.north);

% ---------- (d) tổ chức và lưu trữ
\node[st, minimum height=7mm, text width=14mm] (n1) at (5.45,-5.45) {Môn học};
\node[st, minimum height=7mm, text width=18mm] (n2) at (7.75,-5.45) {Điểm kiến thức};
\node[st, minimum height=7mm, text width=14mm] (n3) at (10.0,-5.45) {Khái niệm};
\draw[arr] (n1) -- (n2);
\draw[arr] (n2) -- (n3);
\draw (n3.south) -- (10.0,-6.0);
\draw (5.35,-6.0) -- (10.0,-6.0);
\foreach \x/\t in {5.35/văn bản, 6.9/ảnh, 8.45/video, 10.0/âm~thanh}{%
  \draw[->] (\x,-6.0) -- (\x,-6.25);
  \node[st, fill=green!6, minimum height=5mm, text width=12.5mm] at (\x,-6.5) {\t};
}
\node[st, text width=60mm, minimum height=8mm] (js) at (7.7,-7.5) {JSON (bản ghi và bảng tham chiếu chéo) $\rightarrow$ triple (11 nhãn quan hệ) $\rightarrow$ RDF (\texttt{rdflib})};
\draw[arr] (7.7,-6.75) -- (js.north);

% ---------- (e) ứng dụng
\node[app] at (1.95,-5.6) {Sinh câu hỏi đa phương thức};
\node[app] at (1.95,-6.3) {Hỏi đáp đa phương thức};
\node[app] at (1.95,-7.0) {Truy hồi chéo phương thức};
\node[app] at (1.95,-7.7) {Trợ lý học tập};

% ---------- nối các ô (c) -> (d) -> (e)
\draw[arr] (11.5,-6.6) -- (10.9,-6.6);
\draw[arr] (4.5,-6.6) -- (3.9,-6.6);
\end{tikzpicture}
